\section{Global bounds and exponent dependence}\label{sec:global}

To apply Proposition~\ref{prop:sparse} globally, we cover each superlevel
set by a controlled number of sparse families. The $L^4$ and gradient
bounds below control the number of occupied unit cubes; slow growth of
local point counts then selects the radii without a diameter bound.

We retain the fixed graph, support, and window from the preceding
sections.  All constants in this section, except for the displayed
dependence on the small exponents, depend only on that fixed setup.
The elementary estimates below use only Equation~\eqref{eq:ellipticity}.
The global estimate for $3<p\leq4$ additionally uses
Proposition~\ref{prop:sparse}, and hence the packet input in
Theorem~\ref{thm:packet}.

\subsection{Elementary bounds and the number of occupied cubes}

\begin{lemma}[Basic graph estimates]\label{lem:basic}
Let $g\in L^2(\mathbb R^2)$ have compact support, and let
\[
 G(y',t)=\int_{\mathbb R^2}g(\xi)
                  e\big(y'\cdot\xi+t\phi(\xi)\big)\,d\xi.
\]
If $\phi''\geq cI$ on $\mathbb R^2$, then
\begin{equation}\label{eq:elementary-lfour}
 \|G\|_{L^4(\mathbb R^3)}
       \leq (\pi/c)^{1/4}\|g\|_{L^2(\mathbb R^2)}.
\end{equation}
In particular, for $F=E_\phi(h\chi_1^{\nu_*})$ with $h$ supported
in $Q_0$, there is a fixed $C_{\mathrm b}\geq1$ such that
\begin{equation}\label{eq:basic-bounds}
 \|F\|_4+\|F\|_\infty+\|\nabla F\|_\infty
       \leq C_{\mathrm b}\|h\|_3.
\end{equation}
\end{lemma}

\begin{proof}
Let $\lambda$ be the finite complex measure on the graph of $\phi$
whose density in base coordinates is $g$.  For
$v\in C_c^\infty(\mathbb R^3)$, Cauchy--Schwarz gives
\begin{align*}
 \left|\int v\,d(\lambda*\lambda)\right|
 &\leq \|g\|_2^2
 \left(\iint_{\mathbb R^2\times\mathbb R^2}
 |v(\xi+\xi',\phi(\xi)+\phi(\xi'))|^2\,d\xi\,d\xi'\right)^{1/2}.
\end{align*}
Put $m=\xi+\xi'$ and $w=(\xi-\xi')/2$; this linear change
has absolute Jacobian one.  Write $w=r\omega$, where
$r>0$ and $\omega\in S^1$, and set
\[
 H_{m,\omega}(r)=\phi(m/2+r\omega)+\phi(m/2-r\omega).
\]
Uniform convexity implies
\[
 H_{m,\omega}'(r)
 =\int_{-r}^{r}\omega^{\mathsf T}\phi''(m/2+s\omega)
                                      \omega\,ds
 \geq2cr.
\]
Thus $H_{m,\omega}$ is strictly increasing for $r>0$, and the
substitution $s=H_{m,\omega}(r)$ yields
\begin{align*}
 &\iint |v(\xi+\xi',\phi(\xi)+\phi(\xi'))|^2\,d\xi\,d\xi'
 \\
 &\quad=\int_{\mathbb R^2}\int_{S^1}\int_0^\infty
        |v(m,H_{m,\omega}(r))|^2r\,dr\,d\omega\,dm
 \leq\frac{\pi}{c}\|v\|_2^2.
\end{align*}
The omitted set $r=0$ has measure zero.  Consequently
$\lambda*\lambda$ has an $L^2$ density with norm at most
$(\pi/c)^{1/2}\|g\|_2^2$, by the $L^2$ representation theorem.
Its Fourier transform is $G(-\,\cdot)^2$ in the sense of
tempered distributions.  Plancherel therefore gives
Equation~\eqref{eq:elementary-lfour}; this reasoning applies to the
stated rough data and does not assume the convolution has a density
in advance.

For $g=h\chi_1^{\nu_*}$, compact support and H\"older's inequality
bound $\|g\|_2$ by a fixed multiple of $\|h\|_3$.
Furthermore
\[
 \|F\|_\infty\leq\|g\|_1,\qquad
 \|\nabla F\|_\infty
 \leq2\pi\int |(\xi,\phi(\xi))|\,|g(\xi)|\,d\xi.
\]
Differentiation under the integral is valid because the integrands
are supported in a fixed compact set.  Another application of
H\"older's inequality proves Equation~\eqref{eq:basic-bounds}.
\end{proof}

For the covering argument, normalize $\|h\|_3\leq1$ and write
\[
 A_\mu=\{y\in\mathbb R^3:|F(y)|>\mu\},\qquad 0<\mu\leq1.
\]
The half-open cubes $k+[0,1)^3$, $k\in\mathbb Z^3$, form a partition.
Call such a cube occupied if it meets $A_\mu$.

\begin{lemma}[Finite occupancy]\label{lem:occupancy}
There is a fixed $C_X\geq1$ such that at most $C_X\mu^{-7}$ unit
cubes are occupied.  In particular, if $\mu\leq C_X^{-1}$, one can
choose one point of $A_\mu$ in each occupied cube to obtain a finite
set $X$ satisfying $\#X\leq\mu^{-8}$.
\end{lemma}

\begin{proof}
First take any finite collection $\mathcal Q$ of occupied cubes
and choose $x_Q\in Q\cap A_\mu$ for each of them.  By
Equation~\eqref{eq:basic-bounds}, on
$B(x_Q,r_\mu)$, where $r_\mu=\mu/(2C_{\mathrm b})\leq1/2$,
one has $|F|>\mu/2$.  These balls have overlap at most $125$:
if a point belongs to one of them, its center lies within distance
one, and at most $5^3$ unit lattice cubes can contain such centers.
Separation of the selected points is not required.  Therefore
\[
 \#\mathcal Q\,\frac{\pi}{96C_{\mathrm b}^3}\mu^7
 \leq\sum_{Q\in\mathcal Q}\int_{B(x_Q,r_\mu)}|F|^4
 \leq125\|F\|_4^4\leq125C_{\mathrm b}^4.
\]
For example, one may take
$C_X=\max\{1,12000C_{\mathrm b}^7/\pi\}$.
This bound holds for every finite collection of occupied cubes.
An infinite collection would have finite subsets of arbitrarily
large cardinality, a contradiction.  Thus the full collection is
finite with the same bound.  Finally
$C_X\mu^{-7}\leq\mu^{-8}$ when $\mu\leq C_X^{-1}$.
\end{proof}

\subsection{A multiscale cover without a diameter bound}

The finite-growth and coloring argument below is the covering step
of the sparse-ball method; compare Bourgain--Guth
\cite[arXiv version~3, Appendix, Lemma~A.2]{BourgainGuth2011}.
We retain explicit scale choices to track the dependence on the output
exponent.

Fix once and for all a number $\mu_0>0$ that is no larger than
the threshold in Proposition~\ref{prop:sparse}, $C_X^{-1}$, and
$1/2$.  In particular, for $0<\mu\leq\mu_0$ we have
\begin{equation}\label{eq:fixed-geometric-thresholds}
 \mu^{-60}\geq8,\qquad \mu^{-100}>2.
\end{equation}
This choice is independent of all small exponents below.

\begin{lemma}[Sparse covering]\label{lem:sparse-cover}
Let $0<\tau\leq1$, $0<\mu\leq\mu_0$, and let $X$ be the set
in Lemma~\ref{lem:occupancy}.  Set
\begin{equation}\label{eq:cover-scales}
 \beta=\frac{\tau}{4},\qquad
 J=\left\lceil\frac{32}{\tau}\right\rceil+1,\qquad
 R_j=\mu^{-100(j+1)}\quad(0\leq j\leq J).
\end{equation}
The union of the occupied cubes is covered by at most
\begin{equation}\label{eq:family-count}
 3J\bigl(1+8\log_2(1/\mu)\bigr)\mu^{-\beta}
\end{equation}
families of balls.  Each family has at most $\mu^{-8}$ balls,
all of radius $R=4R_j$ for some $0\leq j<J$, and its centers
are separated by at least $\mu^{-40}R$.
For the dyadic $N$ chosen in Proposition~\ref{prop:sparse},
\begin{equation}\label{eq:cover-N}
 4R\leq N^2<16R,\qquad N<8\mu^{-50J}.
\end{equation}
\end{lemma}

\begin{proof}
If $X$ is empty there is nothing to cover.  Otherwise, using open
balls throughout, set
\[
 n_j(x)=\#\bigl(X\cap B(x,R_j)\bigr),\qquad x\in X.
\]
For each $x$ there is an index $0\leq j<J$ with
\begin{equation}\label{eq:slow-growth}
 n_{j+1}(x)\leq\mu^{-\beta}n_j(x).
\end{equation}
Indeed, failure at every index would imply
$n_J(x)>\mu^{-\beta J}n_0(x)\geq\mu^{-\beta J}>\mu^{-8}$,
since $\beta J>8$, contrary to $\#X\leq\mu^{-8}$.
Assign each $x$ one such index $j$ and the unique integer $k\geq0$
for which
\[
 2^k\leq n_j(x)<2^{k+1}.
\]
There are at most $J(1+8\log_2(1/\mu))$ nonempty groups
with a common pair $(j,k)$.

In one group, choose a maximal pairwise disjoint collection of
the open balls $B(x,R_j)$, and denote its centers by $Y$.
The collection is finite.  Each unselected ball intersects a
selected one, so every center in the group is within distance
$2R_j$ of a point of $Y$.  A point of the unit cube containing
that center is within a further distance $\sqrt3$.
Since $R_j\geq\mu^{-100}$, the balls
$\{B(y,4R_j):y\in Y\}$ cover all cubes assigned to the group.

Make a finite graph on $Y$, joining distinct $x,y$ when
$|x-y|<R_{j+1}/2$.  For a fixed vertex $x$, the original balls
$B(y,R_j)$ of all its neighbors are disjoint and lie in
$B(x,R_{j+1})$: Equation~\eqref{eq:fixed-geometric-thresholds}
gives $R_j<R_{j+1}/2$.  Each neighbor ball contains at least
$2^k$ points of $X$, whereas
Equation~\eqref{eq:slow-growth} gives
\[
 n_{j+1}(x)<\mu^{-\beta}2^{k+1}.
\]
The degree of $x$ is therefore less than $2\mu^{-\beta}$.
A greedy coloring uses at most $3\mu^{-\beta}$ colors.
Centers of the same color are separated by at least
$R_{j+1}/2$, and
\[
 \frac{R_{j+1}}2
 =\frac{\mu^{-100}}2R_j
 \geq\mu^{-40}(4R_j)
\]
by Equation~\eqref{eq:fixed-geometric-thresholds}.
Each color thus supplies a family satisfying the separation
condition, with at most $\#X\leq\mu^{-8}$ balls and
$R=4R_j\geq\mu^{-100}$.

Multiplying the group and color counts proves
Equation~\eqref{eq:family-count}.  A dyadic $N$ with
$4R\leq N^2<16R$ exists, and satisfies
\[
 N<4\sqrt R=8\mu^{-50(j+1)}\leq8\mu^{-50J}.
\]
The argument uses neither the absolute locations of the points
nor the diameter of $X$.  In particular, different clusters may
be arbitrarily far apart.  The use of open balls also resolves
boundary cases: unselected balls meet selected balls by maximality,
and equality in the graph's distance threshold leaves the desired
weak separation inequality intact.
\end{proof}

\subsection{Integration of the distribution bound}

\begin{theorem}[Global graph estimate]\label{thm:graph}
Fix a phase and window setup for which Theorem~\ref{thm:packet}
and Proposition~\ref{prop:sparse} apply.  For $3<p\leq4$, put
\begin{equation}\label{eq:theta-choice}
 \tau=p-3,\qquad \theta=10^{-6}\tau^2.
\end{equation}
There is a constant $C$, independent of $p$ in this interval,
such that
\begin{equation}\label{eq:global-graph-bound}
 \|E_\phi(h\chi_1^{\nu_*})\|_p
 \leq C\left[
       P(\theta,\theta)M(\theta)^{3/2}\tau^{-7}
       \right]^{1/p}\|h\|_3
\end{equation}
for all $h\in L^3$ supported in $Q_0$.
For $4<p<\infty$, Equation~\eqref{eq:basic-bounds} alone gives
$\|E_\phi(h\chi_1^{\nu_*})\|_p\leq C_{\mathrm b}\|h\|_3$.
\end{theorem}

\begin{proof}
By homogeneity it suffices to take $\|h\|_3\leq1$.
In Proposition~\ref{prop:sparse} choose
$\kappa=\eta=s=\theta$.  These choices lie in the permitted
parameter ranges, and
\[
 \gamma_0=10\kappa+\eta+\tfrac32s=\tfrac{25}{2}\theta.
\]
For $J$ in Equation~\eqref{eq:cover-scales}, we have
\begin{equation}\label{eq:exponent-budget}
 J\leq\frac{34}{\tau},\qquad
 50J\gamma_0\leq\frac{17}{800}\tau<\frac{\tau}{4}.
\end{equation}
Consequently the scale in every family furnished by
Lemma~\ref{lem:sparse-cover} satisfies
\[
 N^{\gamma_0}\leq8^{\gamma_0}\mu^{-\tau/4}
                      \leq2\mu^{-\tau/4},\qquad
 \log(2+N)\leq53J\bigl(1+\log(1/\mu)\bigr).
\]
Applying Proposition~\ref{prop:sparse} to each family and summing
over the cover yields, for $0<\mu\leq\mu_0$,
\begin{equation}\label{eq:global-distribution}
 |A_\mu|\leq
 C P(\theta,\theta)M(\theta)^{3/2}
 J^3\bigl(1+\log(1/\mu)\bigr)^3\mu^{-3-\tau/2}.
\end{equation}
Here the colors cost $\mu^{-\beta}=\mu^{-\tau/4}$, while
the scale factor costs at most $\mu^{-\tau/4}$.
The group count contributes one factor of
$J(1+\log(1/\mu))$, and the logarithm squared in
Proposition~\ref{prop:sparse} contributes the other two.
All constants and $\mu_0$ are independent of $\tau$.

The layer-cake identity, valid by Tonelli even before finiteness
of the norm is known, is
\[
 \|F\|_p^p=p\int_0^\infty\mu^{p-1}|A_\mu|\,d\mu.
\]
For its portion below $\mu_0$, Equation~\eqref{eq:global-distribution}
and $p=3+\tau$ reduce the needed integral to
\begin{align}
 \int_0^1\mu^{-1+\tau/2}\bigl(1+\log(1/\mu)\bigr)^3\,d\mu
 &=\frac2\tau+\frac{12}{\tau^2}
                  +\frac{48}{\tau^3}+\frac{96}{\tau^4}
 \leq158\tau^{-4}.\label{eq:layer-integral}
\end{align}
For example, the substitution $u=\log(1/\mu)$ turns the
integral into $\int_0^\infty e^{-\tau u/2}(1+u)^3\,du$,
which gives the displayed identity by integrating each monomial.
Since $p\leq4$ and $J^3\leq34^3\tau^{-3}$, the small-level
contribution is at most
$C P(\theta,\theta)M(\theta)^{3/2}\tau^{-7}$.

For the remaining levels, again by Tonelli,
\begin{align*}
 p\int_{\mu_0}^\infty\mu^{p-1}|A_\mu|\,d\mu
 &\leq\int_{\{|F|>\mu_0\}}|F|^p
 \leq\mu_0^{p-4}\|F\|_4^4
 \leq\mu_0^{-1}C_{\mathrm b}^4.
\end{align*}
This is uniform for $3<p\leq4$ and is absorbed into the previous
bound because $P,M\geq1$ and $\tau\leq1$.
Taking the $p$th root and restoring homogeneity proves
Equation~\eqref{eq:global-graph-bound}.
If $p>4$, directly use
\[
 \|F\|_p^p\leq\|F\|_\infty^{p-4}\|F\|_4^4
                 \leq C_{\mathrm b}^p\|h\|_3^p.
\]
\end{proof}

The factor $\tau^{-7}$ explicitly records the elementary loss
removal argument.  The assumptions provide finite constants
$P(\theta,\theta)$ and $M(\theta)$, but do not provide rates for
them as $\theta\downarrow0$.  Equation~\eqref{eq:global-graph-bound}
therefore asserts no polynomial upper bound in $(p-3)^{-1}$ for
the full extension constant.

\subsection{A necessary divergence at the endpoint}

\begin{proposition}[Lower bound for the best constant]\label{prop:lower}
For every nonempty compact smooth surface $\Sigma$ in the stated
class there is $c_\Sigma>0$ such that its best extension constant
satisfies
\begin{equation}\label{eq:constant-lower-bound}
 C_{\Sigma,p}^{\mathrm{best}}
       \geq c_\Sigma(p-3)^{-1/p},\qquad 3<p\leq4.
\end{equation}
This conclusion does not use the Kakeya or packet estimates.
\end{proposition}

\begin{proof}
Choose a nonempty interior chart; a nonempty smooth surface with
boundary also has such a chart.  Write it in the affine graph form
\[
 \Psi(\xi)=v_0+L(\xi,\phi(\xi)),\qquad \xi\in D,
\]
where $L$ is invertible, and write $d\sigma=J_\Sigma(\xi)\,d\xi$
on this chart.  Take a fixed nonzero $g\in C_c^\infty(D)$ and
define $f(\Psi(\xi))=g(\xi)/J_\Sigma(\xi)$ on the chart, with
$f=0$ elsewhere.  This is a smooth surface datum with
\[
 B_f:=\sup_{3\leq p\leq4}\|f\|_{L^p(\Sigma)}<\infty.
\]
Let
\[
 G(y',t)=\int g(\xi)e\big(y'\cdot\xi+t\phi(\xi)\big)\,d\xi,
 \qquad A=\|g\|_2^2>0.
\]
Horizontal Plancherel and its derivative identity give, for $t\geq0$,
\begin{align*}
 \int_{\mathbb R^2}|G(y',t)|^2\,dy'&=A,\\
 \int_{\mathbb R^2}|y'|^2|G(y',t)|^2\,dy'
 &=\frac1{(2\pi)^2}
       \left\|\nabla\big(g e(t\phi)\big)\right\|_2^2
 \leq B_g(1+t)^2,
\end{align*}
where one may take
$B_g=2(2\pi)^{-2}\|\nabla g\|_2^2+2\|g\nabla\phi\|_2^2>0$.
Set $K=\max\{1,(2B_g/A)^{1/2}\}$.
The integral outside the disk $|y'|<K(1+t)$ is at most
$B_g/K^2\leq A/2$, so at least $A/2$ of the squared norm lies
inside that disk.  H\"older's inequality on the disk implies
\begin{equation}\label{eq:horizontal-lower-bound}
 \int_{\mathbb R^2}|G(y',t)|^p\,dy'
 \geq (A/2)^{p/2}(\pi K^2)^{1-p/2}(1+t)^{2-p}
 \geq c_g(1+t)^{2-p},
\end{equation}
uniformly for $3\leq p\leq4$, with the explicit positive choice
\[
 c_g=\min\left\{
 (A/2)^{3/2}(\pi K^2)^{-1/2},
 (A/2)^2(\pi K^2)^{-1}
 \right\}.
\]
Integrating Equation~\eqref{eq:horizontal-lower-bound} over
$t\geq0$ yields $\|G\|_p^p\geq c_g/(p-3)$.

The affine coordinates introduce exactly the factors
\[
 E_\Sigma f(x)=e(x\cdot v_0)G(L^{\mathsf T}x),\qquad
 \|E_\Sigma f\|_p^p=|\det L|^{-1}\|G\|_p^p.
\]
Hence, on setting $d_g=|\det L|^{-1}c_g>0$, we obtain
\[
 C_{\Sigma,p}^{\mathrm{best}}
 \geq B_f^{-1}d_g^{1/p}(p-3)^{-1/p}
 \geq B_f^{-1}\min\{d_g^{1/3},d_g^{1/4}\}(p-3)^{-1/p}.
\]
This proves Equation~\eqref{eq:constant-lower-bound}, with every
choice of chart and datum fixed independently of $p$.
\end{proof}
